\question[3]
Ein Verein möchte Nachrichten nur noch mit asymmetrischer Verschlüsselung per email empfangen. Deshalb soll der Webmaster,
 der Zugriff auf den Mailserver hat, den öffentlichen Schlüssel auf der Homepage des Vereins veröffentlichen. 
Erläutere, wie der Webmaster mit Hilfe eines sogenannten \textit{Man-in-the-Middle-Angriffs} die Nachrichten manipulieren kann.

\begin{solutionbox}{4cm}
Der Webmaster veröffentlicht seinen eigenen öffentlichen Schlüssel an Stelle des öffentlichen Schlüssels
des Vereins auf deren Homepage. Wenn ein Kunde mit Hilfe dieses Schlüssels eine Nachricht an den Verein schickt, fängt der Webmaster die Mail auf dem Mailserver ab, entschlüsselt
sie mit Hilfe seines privaten Schlüssels, manipuliert die Nachricht und schickt sie dann mit Hilfe des
öffentlichen Schlüssels des Vereins an den Verein weiter.
\end{solutionbox}
